\begin{titlepage}
\vspace*{1.5cm}
{\Huge\bfseries 高级图像处理\par}
\vspace{.5cm}
{\LARGE 考试复习稿\par}
\vspace{.5cm}
{\large 知识点、公式、例题与逐步解答\par}
\vspace{1.2cm}
\noindent 依据：课程 PPT、Gonzalez 与 Woods《Digital Image Processing》第4版，以及现有五段课堂文字记录。\par
\vspace{.6cm}
\noindent 本稿优先覆盖首次小测前已讲的基础与空间域增强，另列10月12日的频域内容。它是复习讲义，不是老师公布的题库或押题。\par
\vfill
\noindent 资料核对日期：2026年10月19日\par
\noindent 公式采用 LaTeX 排版；自编练习均明确标注，不冒充课堂原题。\par
\end{titlepage}
\section*{先看这一页：考什么，怎样用这份稿}
\addcontentsline{toc}{section}{先看这一页：考什么，怎样用这份稿}
\subsection*{老师明确说过的安排}
\begin{itemize}
\item 10月12日开场：下周用15分钟回答一两个小问题，覆盖前面讲过的内容；强调理解与思考，不强调背诵。按录音日期推算，下周对应10月19日，并非另行取得的考试通知。
\item 全学期三次小测。9月21日说明：参加至少两次算全勤，成绩取最好的一次。
\item 老师明确举过的题型是“简单计算分辨率需求”及简单概念问题。未公布逐章必考清单。
\item 9月28日直方图均衡化部分说理解即可，不再要求推导该部分。\textbf{这不等于不会考映射步骤、结果或理解。}
\item 项目建议书、编程测试、期末报告与展示另成一部分考核；不要把准备小测等同于完成课程考核。
\end{itemize}
\subsection*{范围与优先级}
\textbf{优先：}成像与分辨率、采样和量化、增强与恢复、灰度映射、直方图处理、空间平滑与锐化、滤波器选择。\par
\textbf{补充：}10月12日新讲的频域增强与同态滤波。老师在该课开场说“前面讲过”，无法据此确定当天新内容是否也纳入第一次小测，因此单独标出。\par
几何校正、压缩等后续课程内容，不因教材中出现就擅自列为本次小测重点。开闭卷、成绩百分比与具体试题尚未确认。
\subsection*{建议使用顺序}
先读每节概念和易错点，再遮住答案做例题。时间紧张时，优先做分辨率需求、灰度映射、直方图均衡化和邻域滤波的手算。最后用本稿末尾的15分钟练习检查自己能否讲清原因。
\subsection*{资料使用说明}
课堂文字记录可能存在转写误差。本稿用教材核对数学定义，并在课堂数值或术语不明确时说明假设。所有补充题均为教学自编，数值结果可以独立复算。部分公式是帮助理解课堂概念的标准表达，不意味着老师要求默写。
\clearpage
\tableofcontents
\clearpage
